\documentclass[10pt,twocolumn]{article}
\usepackage[margin=0.75in]{geometry}
\usepackage{times}
\usepackage{mathptmx}
\usepackage{amsmath,amssymb,amsfonts}
\usepackage{booktabs}
\usepackage{graphicx}
\usepackage{microtype}
\usepackage{natbib}
\usepackage{hyperref}
\usepackage{url}
\usepackage{enumitem}
\usepackage{caption}
\usepackage{subcaption}

\setlength{\columnsep}{0.24in}
\setlength{\textfloatsep}{8pt plus 2pt minus 2pt}
\setlength{\floatsep}{6pt plus 2pt minus 2pt}
\setlength{\intextsep}{6pt plus 2pt minus 2pt}
\setlength{\abovecaptionskip}{4pt}
\setlength{\belowcaptionskip}{0pt}
\setlist[itemize]{leftmargin=1.1em,topsep=2pt,itemsep=1pt,parsep=0pt}
\setlist[enumerate]{leftmargin=1.2em,topsep=2pt,itemsep=1pt,parsep=0pt}
\hypersetup{colorlinks=true,linkcolor=blue,citecolor=blue,urlcolor=blue}

\title{Angular Manifold Routing: Sublinear Compute Reduction via Hopf-Base Sector Discretization}
\author{Casey Allard\\Independent Researcher}
\date{}

\begin{document}
\maketitle

\begin{abstract}
We study whether the angular non-uniformity of normalized token embeddings can be exploited for routing computation in transformer-style architectures. We propose a fixed geometric routing method based on a Hopf-fibration-derived sector map and evaluate it in two settings: a semantic proxy based on PPMI-SVD embeddings and a small trainable language model. On the semantic proxy, the effective routing footprint grows sublinearly with routing budget, following an empirical scaling law of approximately $K^{0.572}$, compared with $K^{1.0}$ for dense routing, with the advantage persisting through $K=5000$. In a 2-layer trainable language model, fixed geometric routing replaces a learned top-1 gate with an 8\% validation perplexity cost and no learned gate matrix, while retaining a concentrated routing footprint. A second-dataset replication on WikiText-2 matches the Penn Treebank result closely. A key limitation is that in the trainable setting, the specific Hopf sector geometry is not distinguishable from randomly permuted fixed routing, suggesting that expert weights co-adapt to any fixed routing scaffold in this small-scale regime. The contribution of this work is therefore narrow: angular structure supports sparse fixed routing and transfers partially to trainable routing, but geometry-specific quality gains are not established in end-to-end language model training.
\end{abstract}

\section{Introduction}
Sparse routing is a standard strategy for reducing the effective compute cost of transformer-style models. Mixture-of-experts (MoE) architectures activate only a subset of available compute paths for a given input, typically using a learned gating network \citep{shazeer2017outrageously,fedus2022switch}. This paper studies a different question: can some routing structure be obtained directly from the geometry of the representation space, without a learned routing network?

The motivating observation is that normalized token embeddings are not uniformly distributed on the unit hypersphere. Recent compression work has shown that this angular structure can be exploited for efficient representation \citep{turboquant2026}. Here we ask whether the same broader geometric fact can support fixed routing. Our approach uses a fixed routing map derived from a 4D Hopf-fibration projection.

This paper makes three scoped contributions. First, on a semantic proxy, fixed geometric routing yields a sublinear effective routing footprint over a broad range of routing budgets. Second, in a small trainable language model, fixed routing partially transfers, replacing a learned top-1 gate at modest perplexity cost. Third, in that trainable setting, the specific Hopf geometry does not outperform randomly permuted fixed routing, clarifying the scope of the result.

The intended claim is narrow. This work does \emph{not} establish large-scale replacement of learned MoE routing in realistic systems. It presents an empirical efficiency result in a proxy setting and a toy-scale trainable language model.

\section{Background}
Compression work such as TurboQuant treats angular structure in normalized embeddings as a representation property to be transformed for scalar quantization \citep{turboquant2026}. Our focus is different: we ask whether a fixed geometric partition of the sphere can induce a non-uniform routing distribution that concentrates traffic into a subset of available sectors.

If semantically related points cluster into preferred angular regions, then a sector map on the sphere should receive non-uniform traffic. In principle, this can reduce the effective routing footprint without introducing a learned gate matrix.

\section{Method}
\subsection{Fixed geometric routing}
Given an L2-normalized input vector $\mathbf{x} \in \mathbb{R}^D$, we project to a 4D subspace using four fixed coordinates,
\[
(a,b,c,d) = (x_i,x_j,x_k,x_l),
\]
with canonical index choice $(3,65,2,21)$ in the reported experiments.

We then compute Hopf-base coordinates
\begin{align}
\rho_1 &= \sqrt{a^2+b^2}, & \rho_2 &= \sqrt{c^2+d^2}, \\
\chi_u &= \frac{\rho_2^2}{\rho_1^2+\rho_2^2},
& \delta &= \operatorname{atan2}(b,a)-\operatorname{atan2}(d,c), \\
\alpha &= \tfrac{1}{2}\left(\operatorname{atan2}(b,a)+\operatorname{atan2}(d,c)\right).
\end{align}
An optional transport correction is applied as
\[
\alpha_t = \alpha + \frac{1}{2}\lambda \cos(2\chi)\,\delta,
\]
where $\chi$ is derived from $(\rho_1,\rho_2)$ and $\lambda$ controls the correction strength. The resulting coordinates are discretized into a product sector budget $K$, yielding a fixed routing sector.

\subsection{Evaluation metric}
To quantify routing concentration, we use an effective-buckets metric $\hat e$ that summarizes how many routing sectors are effectively active under the empirical routing distribution. Lower $\hat e$ relative to $K$ indicates stronger concentration and therefore a smaller effective routing footprint.

\section{Proxy Experiments}
\subsection{Setup}
We first evaluate the routing mechanism on a semantic proxy based on PPMI-SVD embeddings derived from standard distributional-semantics pipelines \citep{mikolov2013word2vec,pennington2014glove}. This setting isolates routing geometry from end-to-end expert co-adaptation.

The reported proxy study spans more than 160 runs over routing budgets from $K=25$ to $K=5000$. Comparisons include the canonical Hopf configuration, a randomized fixed-routing control, and dense routing.

\subsection{Main results}
Across the proxy experiments, the effective routing footprint for the canonical Hopf configuration is well fit by
\[
\hat e \approx 2.957 \cdot K^{0.572},
\]
compared with linear $K^{1.0}$ scaling for dense routing. A randomized control scales as approximately $1.664\cdot K^{0.814}$, intermediate between Hopf and dense routing.

\begin{table}[t]
\centering
\caption{Proxy routing scaling laws.}
\label{tab:scaling}
\small
\begin{tabular}{lcc}
\toprule
Routing & Scaling law & $R^2$ \\
\midrule
HOPF & $\hat e = 2.957 \cdot K^{0.572}$ & 0.963 \\
PERMUTED & $\hat e = 1.664 \cdot K^{0.814}$ & 0.993 \\
DENSE & $\hat e = K^{1.0}$ & 1.000 \\
\bottomrule
\end{tabular}
\end{table}

\begin{figure}[t]
    \centering
    \includegraphics[width=\columnwidth]{figures/fig_scaling_law.pdf}
    \caption{Effective routing footprint versus routing budget. Fixed Hopf routing shows sublinear growth relative to dense routing.}
    \label{fig:scaling}
\end{figure}

The advantage persists through $K=5000$, where the observed ratio relative to dense routing remains in the range of approximately $2.6$--$2.8\times$. The corresponding hardware-oriented proxy analysis reports $3.0$--$4.9\times$ lower effective cost and $2.5$--$2.9\times$ fewer LRU-16 cache misses than dense routing. The scaling exponent is also approximately norm-invariant, varying by less than $0.015$ across L1/L2/L3/L4 normalization schemes.

\subsection{Mechanism}
An ablation on the proxy indicates that the dominant signal comes from the raw fiber phase $\alpha=(\theta_1+\theta_2)/2$, while the transport correction contributes a smaller additional gain. In the reported proxy measurements, the raw phase term accounts for the majority of the concentration advantage, with the transport term adding a modest improvement on top.

\section{Trainable Language Model Experiments}
\subsection{Setup}
We evaluate the method in a 2-layer transformer language model with hidden size 128, four attention heads, sequence length 32, vocabulary size 5000, and routing budget $K=64$. Three conditions are compared:
\begin{itemize}
    \item \textbf{BASELINE}: learned top-1 routing gate,
    \item \textbf{HOPF-ROUTED}: fixed geometric routing with no learned gate matrix,
    \item \textbf{PERMUTED}: randomly permuted fixed routing sectors.
\end{itemize}
All conditions are approximately parameter matched at $\ge 5.6$M parameters. Training uses 4000 steps with cosine learning rate $3\times10^{-4}$ and batch size 64. Evaluation is reported on Penn Treebank \citep{ptb1993} and WikiText-2 \citep{merity2016wt2}, each confirmed across two seeds.

\subsection{Penn Treebank}
On Penn Treebank, fixed geometric routing reaches within 8\% of the learned-gating baseline in validation perplexity while retaining a concentrated routing footprint and using no learned gate matrix. Table~\ref{tab:ptb} reports the two-seed mean result.

\begin{table}[t]
\centering
\caption{Penn Treebank results (2-seed mean, 4000 steps). Lower effective-path count relative to $K=64$ indicates stronger concentration.}
\label{tab:ptb}
\small
\begin{tabular}{lccc}
\toprule
Method & Val PPL & Eff. paths & Eff. ratio vs. dense \\
\midrule
BASELINE & 152.26 & 43.19 & $1.48\times$ \\
HOPF-ROUTED & 164.54 & 45.96 & $1.39\times$ \\
PERMUTED & 164.41 & 45.81 & $1.40\times$ \\
\bottomrule
\end{tabular}
\end{table}

The HOPF/BASELINE validation-perplexity ratio is $1.081$. HOPF-ROUTED and PERMUTED differ by only $0.13$ perplexity.

\subsection{WikiText-2 replication}
The same comparison on WikiText-2 yields a nearly identical HOPF/BASELINE ratio under the same training setup. Table~\ref{tab:wt2} shows the per-seed detail, and Table~\ref{tab:crossdataset} summarizes the cross-dataset comparison.

\begin{table}[t]
\centering
\caption{WikiText-2 replication (2-seed detail).}
\label{tab:wt2}
\small
\begin{tabular}{lccc}
\toprule
Method & Seed 0 PPL & Seed 1 PPL & Mean PPL \\
\midrule
BASELINE & 106.96 & 104.74 & 105.85 \\
HOPF-ROUTED & 113.65 & 115.25 & 114.45 \\
PERMUTED & 113.86 & 115.09 & 114.48 \\
\bottomrule
\end{tabular}
\end{table}

\begin{table}[t]
\centering
\caption{Cross-dataset summary.}
\label{tab:crossdataset}
\small
\begin{tabular}{lcc}
\toprule
Metric & PTB & WikiText-2 \\
\midrule
HOPF/BASELINE PPL ratio & 1.081 & 1.081 \\
HOPF vs. PERMUTED $|\Delta \mathrm{ppl}|$ & 0.13 & 0.03 \\
HOPF vs. DENSE effective-path ratio & $1.39\times$ & $1.56\times$ \\
\bottomrule
\end{tabular}
\end{table}

\subsection{Interpretation}
The trainable experiments support a narrow transfer claim: fixed geometric routing can replace a learned top-1 gate at moderate validation-perplexity cost while preserving a concentrated routing footprint. At the same time, the experiments identify an equally important null result: in the trainable setting, the specific Hopf geometry is not distinguishable from randomly permuted fixed routing. This suggests that expert weights co-adapt to the fixed routing scaffold, absorbing geometry-specific differences in this small-scale regime.

A related control supports the choice of baseline. The reported learned-routing baseline uses top-1 gating without an auxiliary load-balancing loss. In internal control experiments, imposing explicit load-balance regularization drove routing close to uniform and slightly worsened perplexity, indicating that \emph{native learned concentration} is the more appropriate comparison for this study.

\section{Discussion}
The empirical picture is therefore mixed but informative. On the one hand, the proxy experiments show a strong sublinear routing footprint derived from embedding geometry alone. On the other hand, end-to-end language-model training absorbs the geometry-specific advantage, leaving a narrower result: fixed routing can replace learned gating at moderate cost in a toy-scale setting, but the particular geometric scaffold is not uniquely beneficial in that trainable regime.

This positions the contribution as a narrow efficiency result rather than a broad architecture claim. The proxy experiments establish a geometry-derived routing footprint. The trainable experiments establish partial transfer of that footprint and simultaneously identify a null result: the specific Hopf sector geometry does not appear to matter once expert weights co-adapt.

\section{Limitations}
This work has several important limitations.
\begin{enumerate}
    \item The strongest scaling-law results are established on a semantic proxy rather than large end-to-end language models.
    \item The trainable experiments use a 2-layer model and should not be interpreted as evidence for large-scale replacement of learned routing.
    \item The Hopf geometry does not outperform randomly permuted fixed routing in the trainable language model setting.
    \item The experiments do not address attention-side routing.
    \item All experiments were conducted by a single researcher, and external replication is not yet available.
\end{enumerate}

\section{Reproducibility and Conclusion}
A standalone reproduction bundle, including a NumPy-only implementation of the routing function, figure-generation code, and manuscript source, is available at \url{https://doi.org/10.5281/zenodo.19243034}.

We presented a fixed geometric routing method based on a Hopf-fibration-derived sector map and evaluated it on both a semantic proxy and a small trainable language model. The main empirical result is that angular structure in normalized embeddings supports a sublinear effective routing footprint and can partially transfer to a trainable routing setting without a learned gate matrix. The main limitation is that the specific geometry does not provide an observable quality advantage over random fixed routing in the trainable regime studied here.

The contribution is therefore narrow but concrete: embedding angular structure appears relevant not only for compression but also for routing computation.

\begin{thebibliography}{9}

\bibitem[Google Research(2026)]{turboquant2026}
Google Research. 2026.
\newblock {TurboQuant}: Extreme {KV}-Cache Compression via Angular Rotation-Invariant Quantization.
\newblock \emph{arXiv preprint}.

\bibitem[Vaswani et~al.(2017)]{vaswani2017attention}
Ashish Vaswani, Noam Shazeer, Niki Parmar, Jakob Uszkoreit, Llion Jones, Aidan N. Gomez, Lukasz Kaiser, and Illia Polosukhin. 2017.
\newblock Attention Is All You Need.
\newblock \emph{arXiv:1706.03762}.

\bibitem[Shazeer et~al.(2017)]{shazeer2017outrageously}
Noam Shazeer, Azalia Mirhoseini, Krzysztof Maziarz, Andy Davis, Quoc Le, Geoffrey Hinton, and Jeff Dean. 2017.
\newblock Outrageously Large Neural Networks: The Sparsely-Gated Mixture-of-Experts Layer.
\newblock In \emph{Proceedings of ICLR}.

\bibitem[Fedus et~al.(2022)]{fedus2022switch}
William Fedus, Barret Zoph, and Noam Shazeer. 2022.
\newblock Switch Transformers: Scaling to Trillion Parameter Models with Simple and Efficient Sparsity.
\newblock \emph{Journal of Machine Learning Research}, 23:1--39.

\bibitem[Mikolov et~al.(2013)]{mikolov2013word2vec}
Tomas Mikolov, Ilya Sutskever, Kai Chen, Greg S. Corrado, and Jeff Dean. 2013.
\newblock Distributed Representations of Words and Phrases and their Compositionality.
\newblock In \emph{Advances in Neural Information Processing Systems}.

\bibitem[Pennington et~al.(2014)]{pennington2014glove}
Jeffrey Pennington, Richard Socher, and Christopher D. Manning. 2014.
\newblock {GloVe}: Global Vectors for Word Representation.
\newblock In \emph{Proceedings of EMNLP}.

\bibitem[Marcus et~al.(1993)]{ptb1993}
Mitchell P. Marcus, Mary Ann Marcinkiewicz, and Beatrice Santorini. 1993.
\newblock Building a Large Annotated Corpus of English: The Penn Treebank.
\newblock \emph{Computational Linguistics}, 19(2):313--330.

\bibitem[Merity et~al.(2016)]{merity2016wt2}
Stephen Merity, Caiming Xiong, James Bradbury, and Richard Socher. 2016.
\newblock Pointer Sentinel Mixture Models.
\newblock \emph{arXiv:1609.07843}. Introduces WikiText-2 and WikiText-103 benchmarks.

\end{thebibliography}

\end{document}
